\section*{Caso, contrato y alcance}
\textbf{Responsable y turno:} Khriz Alexander Coronel Gómez, minutos 3:00 a 6:00. El caso es registrar una solicitud individual de refuerzo académico en SGRA, conforme al RF01 del PFC~\cite{sgra-rf01-khriz}.

El contexto documentado del proyecto describe acuerdos entre estudiantes y docentes por WhatsApp, correo y encuentros virtuales. El SGRA anterior implementa un formulario y un backend para registrar solicitudes. Como cliente futuro se plantea una aplicación externa de estudiantes que enviaría los mismos datos a una API. Esta demo representa ese intercambio con datos ficticios.

\textbf{Transparencia de acceso:} el cliente utiliza una operación HTTP sin conocer las tablas ni su forma de almacenamiento. En la futura distribución, la misma interfaz podría ocultar qué nodo atiende. \textbf{Límite visible:} una respuesta ausente no permite saber si el registro se guardó; se debe informar que el resultado no está confirmado.

\textbf{Contrato simplificado:} \texttt{POST /api/student/requests}, con JSON. Se requieren \texttt{subjectId} entero positivo, \texttt{sessionTypeId = 1}, \texttt{reason} con al menos diez caracteres tras retirar espacios exteriores y \texttt{preferredSchedule} no vacío. El horario expresa disponibilidad del estudiante; no confirma una sesión.

Con datos válidos, se crea una solicitud en memoria con estudiante ficticio 1, identificador y estado Pendiente. Se devuelve \texttt{201 Created} con \texttt{requestId}. Sin motivo se devuelve \texttt{400}, se identifica el campo y no se inserta. El JSON mal formado también recibe \texttt{400}. Los archivos de peticiones, respuestas y resultados se conservan en \texttt{demos/khriz/}.

El backend anterior usa \texttt{multipart/form-data}, devuelve \texttt{200} y no incluye \texttt{preferredSchedule} en el DTO consultado. El prototipo adapta ese caso a la ficha de EXP-1 y al RF01 reformulado. No autentica a estudiantes ni valida matrícula, docente, período, existencia de asignatura, adjuntos o grupos. Los registros se pierden al reiniciar. No se modifica el PFC ni se presenta esta demo como una integración real.

\textbf{Elección y alternativa:} el intercambio HTTP permite solicitar la creación y recibir inmediatamente su resultado. Una alternativa sería admitir un mensaje en una cola y consultar posteriormente su procesamiento. Kazanavi{\v{c}}ius y Ma{\v{z}}eika comparan comunicación síncrona y asíncrona; la sección 2, p.~3, explica que el cliente síncrono espera la respuesta y depende de la disponibilidad del servicio~\cite{kazanavicius-comunicacion-2023-khriz}. Esa dependencia es un costo del escenario. No se extrapolan al SGRA las mediciones de rendimiento del artículo.

\section*{Respuesta perdida e idempotencia propuesta}
Si la API guarda una solicitud pero se pierde la respuesta, el cliente desconoce si la operación terminó. Reenviar el POST puede crear un duplicado. Esto se observa en el prototipo: repetir una petición válida genera otro identificador.

La propuesta utiliza una clave por operación, conservada durante sus reintentos y asociada al solicitante, al contenido y al resultado. El servidor guardaría el registro y esa relación atómicamente en almacenamiento duradero; una repetición de la misma operación recuperaría el identificador original. Una clave reutilizada con contenido distinto se rechazaría. Una solicitud nueva tendría otra clave. Este mecanismo queda propuesto, no implementado ni probado en la demo.

\newpage
\section*{CAP y PACELC en una evolución futura}
El diagrama distingue la creación local de una solicitud pendiente de una futura confirmación de cupos. La segunda operación supone dos nodos responsables del mismo cupo. No basta con tener dos procesos para demostrar el compromiso de CAP.

Durante una partición, si el servicio no puede coordinar la decisión con la autoridad necesaria para garantizar un único cupo, se pospondría o rechazaría su confirmación. Así se preservaría esa regla a costa de disponibilidad para esa operación~\cite{gilbert-lynch-cap-khriz}. No se afirma que una falla de enlace obligue a detener todas las funciones ni que un timeout HTTP pruebe una partición.

Sin partición, PACELC permite explicar el compromiso entre latencia y consistencia: esperar coordinación para una confirmación consistente añade tiempo al camino de la operación~\cite{abadi-pacelc-khriz}. No se aportan mediciones del SGRA. La idempotencia evita repetir una misma operación; por sí sola no coordina dos confirmaciones distintas del último cupo. REST tampoco determina la prioridad CAP de un sistema.

\section*{Guion de tres minutos y demostración}
\begin{tabularx}{\textwidth}{@{}p{1.8cm}X@{}}
\toprule
Tiempo & Contenido que explica Khriz \\
\midrule
0:00--1:00 & Solicitud de refuerzo, actores, ejemplo conocido, transparencia de acceso y límite de una respuesta ausente. \\
1:00--1:45 & POST válido: 201 e identificador. Sin motivo: 400 y ningún registro nuevo. Mostrar respuestas conservadas o ejecutar las dos peticiones. \\
1:45--2:15 & Registro guardado y respuesta perdida. Riesgo de duplicados y clave de idempotencia propuesta. \\
2:15--3:00 & Dos nodos y último cupo: decisión ante partición, latencia de coordinación y distinción frente al prototipo. \\
\bottomrule
\end{tabularx}

\textbf{Preparación:} desde \texttt{demos/khriz/}, ejecutar \texttt{node servidor-rest.mjs}. En otra terminal, enviar los dos comandos \texttt{curl.exe} del README de la demo. Para comprobar los casos automáticamente, ejecutar \texttt{node probar-rest.mjs}; crea y cierra su propio servicio en un puerto libre. El parámetro \texttt{--guardar-evidencias} conserva peticiones, respuestas y resultados de esa ejecución. La instancia manual escucha en \texttt{127.0.0.1:8083}; se detiene con Ctrl+C.

\textbf{Preguntas para practicar:} ¿Por qué 201 y no 200 en esta demo? ¿Por qué la disponibilidad preferida no es una sesión confirmada? ¿Qué ocurre al repetir el POST? ¿Qué diferencia hay entre idempotencia y garantizar un cupo único? ¿Qué decisión CAP sería distinta para una consulta informativa?

\section*{Aportes, asistencia y cierre}
\begin{tabularx}{\textwidth}{@{}p{2.1cm}X@{}}
\toprule
Responsable & Material y participación \\
\midrule
Khriz & Caso del SGRA y solicitud de implementación. Debe revisar, comprender y defender el material, aportar el bosquejo auténtico y conservar los enlaces a sus commits reales para la tabla grupal. \\
Codex & Apoyo en implementación, pruebas, diagramas, redacción, búsqueda y comprobación de fuentes. \\
\bottomrule
\end{tabularx}
La captura del artículo se prepara con asistencia de Codex; no acredita una lectura personal de Khriz. Quedan a cargo del equipo el ensayo con rostro y voz, su enlace abierto, la fotografía del bosquejo con nombres y fecha y la confirmación docente del enfoque de los cuatro PFC. Khriz coordina el ensayo y prepara una pregunta crítica para otro equipo. No se atribuyen reuniones, commits, asistencia o mediciones sin evidencia.

\newpage
\section*{Fuentes y localizadores}
El artículo de 2023 se consulta en su edición editorial, pp.~1--36, DOI \texttt{10.31577/cai\_2023\_1\_1}. La captura anotada corresponde a la p.~3 y al apartado 2, sobre comunicación síncrona. La página oficial de la revista declara indexación en Scopus: \url{https://www.cai.sk/ojs/index.php/cai/index}. Esto se registra como declaración editorial; no equivale a una consulta independiente del índice.

Las publicaciones de CAP de 2002 y PACELC de 2012 son fuentes fundacionales complementarias. No se cuentan dentro del mínimo de referencias de los últimos cinco años. El equipo debe completar y verificar sus seis referencias recientes y sus tres artículos de revista indexada.

\printbibliography[heading=none]
\section*{Anexos pendientes de integración grupal}
\begin{itemize}
\item Fotografía auténtica del bosquejo a mano, con nombres y fecha. No se sustituye por el diagrama digital.
\item Tres enlaces a commits propios significativos en las fechas reales y la tabla de aportes del grupo.
\item Enlace del ensayo accesible en ventana privada. Este documento individual se integra y ajusta con las demás secciones al límite grupal de 2--4 páginas.
\end{itemize}
